% !TEX root = ../main.tex
% =====================================================================
% APPENDIX B -- REPRODUCIBILITY AND DETAILS     target: 1 - 1.5 pages
%
% The place for the details the examiner did not want in the main text:
% the full hyperparameter list, the exact commands, the environment.
% Short and factual.
% =====================================================================

\section{Reproducibility and Implementation Details}
\label{app:repro}

\begin{guide}
Three short subsections: hyperparameters, commands, environment. No prose needed
beyond one sentence per subsection.
\end{guide}

\subsection{Full hyperparameters}
\label{app:hparams}

\todo{One sentence: every model in this report is defined by one YAML file, and the
table below lists the values shared by all of them.}

{Every model in this report is one YAML file under \texttt{configs/models/}
(\Cref{sec:setup:pipeline}); \Cref{tab:hparams} lists the training recipe and pooling
defaults that are identical across every config I trained, and separates out the one
group that is \emph{not} shared -- encoder width and depth, which differ between the
wide base encoder and the smaller sv4-family variants used for the parameter-count
comparison in \Cref{tab:cost} -- rather than presenting a single number that would be
wrong for half the sweep.}

\begin{table}[h]
\caption{Hyperparameters shared by every model config, plus the two encoder sizes
actually used (see the note below the table for why encoder width/depth is not a
single shared value).}
\label{tab:hparams}
\begin{center}
\small
\begin{tabular}{p{0.30\linewidth} p{0.28\linewidth} p{0.32\linewidth}}
\toprule
\textbf{Group} & \textbf{Parameter} & \textbf{Value} \\
\midrule
Encoder (wide, e.g.\ \texttt{sv2/seanet.yaml}) & width $\dmodel$        & 128 \\
        & number of blocks             & 6 \\
Encoder (narrow, sv4 family, e.g.\ \texttt{seanet\_bottleneck\_topk}) & width $\dmodel$ & 64 \\
        & number of blocks             & 4 \\
\emph{Shared by both} & kernel sizes $\mathcal{K}$   & $\{5, 11, 23\}$, summed per block \\
        & dilation cap $\delta_{\max}$ & 16 \\
        & dropout                      & 0.2 \\
\midrule
Pooling & attention width $\dattn$     & 8 \\
        & dropout                      & 0.2 \\
        & positional encoding          & on \\
        & Top-$k$ fraction $\kappa$ (\Cref{sec:pool:topk}) & 0.1 \\
        & other head-specific ($\beta_0$, $\tau_0$, $\lambda_0$) & \tbd \\
\midrule
Training& optimiser, learning rate    & Adam, $1.25\!\times\!10^{-3}$ \\
        & weight decay                 & $1.2\!\times\!10^{-4}$ \\
        & batch size                   & $\mathrm{clamp}(n_{\text{train}}/10, 2, 16)$ \\
        & epochs, patience             & 400, 60 \\
        & label smoothing              & 0.13 \\
        & focus penalty $\lambda_{\text{entropy}}$ (\Cref{eq:loss}) & 0.01 \\
        & seeds                        & 1 (66-model sweep), 3 for the 3
          \webtraffic{} models in \Cref{tab:seeds} \\
\bottomrule
\end{tabular}
\end{center}
\end{table}
% Source: configs/models/sv2/seanet.yaml (wide encoder, d=128/n_blocks=6, training recipe),
% configs/models/sv4/seanet_bottleneck.yaml and sv4/seanet_inputgate.yaml (narrow encoder,
% d=64/n_blocks=4 -- both share this width, only the block type differs), configs/models/sv6/
% seanet_bottleneck_mschan.yaml (top_frac: 0.1, the only head-specific pooling parameter I could
% find explicitly set in any config file). Kernels/dilation cap/dropout/training recipe are
% identical across every config file I opened (sv1, sv2, sv4, sv6) -- copy-pasted deliberately, per
% each file's own header comment, so results only change encoder-by-encoder or pooling-by-pooling.
% I could NOT find beta_0 / tau_0 / lambda_0 (the Adaptive class-wise, Softmax and Gated-attention
% heads' own internal constants) exposed in any YAML -- if they are tunable, they must be hardcoded
% in seanet/pooling.py, which I did not have open in this session; left \tbd rather than guess.
% lambda_entropy is a TRAINING recipe weight (multiplies the focus penalty in the loss), not a
% pooling-head parameter -- listed under Training here to match \Cref{eq:loss} in the main text.

\subsection{Commands}
\label{app:commands}

\todo{One sentence: the whole study is reproduced by the commands below, run at the
project root.}

{The whole study -- every model, every dataset, every table and figure in this
report -- is reproduced by the three commands below, run from the project root
against one of the config files under \texttt{configs/models/}
(\Cref{sec:setup:pipeline}); no step needs a hand edit to any code file.}

\begin{verbatim}
# train one model over WebTraffic + the UCR archive
python main.py train --model <config>

# comparison tables and the cross-model ranking
python main.py results

# every figure and LaTeX table used in this report
python main.py paper
\end{verbatim}

\todo{One sentence: figures and tables in this report are produced by the last
command, so no number was typed by hand.}

{Every figure and table in this report is produced by \texttt{python main.py paper}
directly from the CSVs \texttt{python main.py train}/\texttt{results} already wrote,
so no number in this report was typed in by hand from a plot or a printed line --
which is also why the staleness issue flagged repeatedly in \Cref{sec:results} and
\Cref{app:additional} is fixable by re-running this one command, not by editing any
table by hand.}

\subsection{Software and hardware environment}
\label{app:env}

\todo{One sentence: development happened locally, the actual training sweep ran on
Grid'5000, and the two environments differ.}

{Day-to-day development (writing and unit-testing the encoder and pooling code,
\Cref{sec:skills}) happened on a local Windows machine with a 6\,GB RTX GPU, used for
writing code and small smoke tests (\Cref{sec:setup}); the 66-model sweep itself ran
on Grid'5000 \citep{balouek2013grid5000}, split across two sites so a delay on one did
not stall the whole study (\Cref{sec:setup:pipeline}), which is why
\Cref{tab:environment} gives the two environments separately rather than as one row.}

\begin{table}[h]
\caption{Environment. Grid'5000 columns are what actually produced every number in
this report; a few local-machine and Grid'5000 fields are left \tbd where I do not
have a way to state them precisely from inside this session (\Cref{app:hparams}'s
introduction explains why I avoid guessing at numbers I cannot check).}
\label{tab:environment}
\begin{center}
\small
\begin{tabular}{lll}
\toprule
 & \textbf{Local (development)} & \textbf{Grid'5000 (training)} \\
\midrule
OS              & Windows & \tbd (site-managed) \\
Python          & \tbd & \tbd \\
Conda env       & --   & \texttt{seanet} (same env on both sites) \\
PyTorch         & \tbd & 2.0.1 \\
CUDA            & \tbd & \tbd \\
GPU              & 1$\times$ RTX (6\,GB), smoke tests only & \tbd \\
Sites / clusters & --  & Lille (\texttt{chuc}, \texttt{chifflot}); Sophia
                        (\texttt{esterel40}, \texttt{esterel43}) \\
Job scheduler   & --   & OAR \\
\bottomrule
\end{tabular}
\end{center}
\end{table}
% GPU/OS row for the Local column: Section 7 (02_context.tex) already states this in prose --
% "a local Windows machine with a 6GB RTX GPU" -- used for writing code and small smoke tests, NOT
% for the actual training sweep, which is why the Grid'5000 columns are the ones that matter for
% reproducing any number in this report. Exact RTX model (e.g. 3060 vs 4060) and exact Windows
% version were not given in that sentence, so I did not invent them.
% Source: requirements_versioned.txt (torch==2.0.1, plus mlflow==3.14.0, optuna==4.9.0,
% pandas==2.0.3, numpy==1.24.4, scikit-learn==1.3.0 -- not all in the table above since it only
% asks for OS/Python/PyTorch/CUDA/GPU/scheduler, but worth keeping in this comment for anyone
% checking exact package versions); GRID5K_CMD_HELP.md (conda env name "seanet", `-p chuc` default
% Lille cluster with `chifflot` given as an alternative, `-p "cluster='esterel40'"`/`'esterel43'`
% for Sophia, OAR via oarsub). I could NOT find a specific GPU model (V100/A100/etc.) anywhere in
% GRID5K_CMD_HELP.md or the configs I had open -- the guide only says "whatever node names appear
% [in the GPU list] have GPUs" without naming the card, so I left GPU \tbd on both sides rather than
% guess; fill it in from the Grid'5000 hardware pages (site-specific, e.g. the Lille or Sophia
% hardware listing) or from `nvidia-smi` output you may already have saved from a session. Local
% (Windows) OS/Python/CUDA are genuinely unknowable to me in this session -- fill from your own
% machine.

\subsection{Code availability}
\label{app:code}

\todo{Two sentences: where the code lives (repository, branch), how it is organised
(one module per stage of the pipeline), and the licence / access conditions if there
are any. Check with my supervisor before publishing a public link.}

{The codebase is organised as one module per stage of the pipeline under
\texttt{seanet/} -- data loading, the encoder and pooling implementations
(\Cref{sec:architecture}, \Cref{sec:pooling}), training, and results/paper generation
(\Cref{sec:setup:pipeline}) -- driven by the config files under \texttt{configs/}
(\Cref{app:hparams}). It is released under the Apache License 2.0 (see
\texttt{LICENSE}) and, per \texttt{NOTICE}, is built on top of Amazon's
\texttt{MILTimeSeriesClassification} codebase (Copyright Amazon.com, Inc.\ or its
affiliates), which is the original implementation of \millet{}
\citep{early2024millet}, so any redistribution needs to keep that attribution and the
licence terms it came with.}

\note{I have deliberately not written a specific repository URL, branch name or public
release date here, per the original \texttt{\textbackslash todo}'s own instruction to
check with my supervisor first -- I do not have a way to confirm that decision from
inside this session, and getting it wrong (e.g.\ publishing a link that was not meant
to be public yet) is not a mistake I want to make on your behalf. What I \emph{can}
confirm from the files themselves: the licence is Apache 2.0, and the codebase carries
Amazon's copyright notice for the upstream MILLET code it extends -- both stated
above. Once you have confirmed the access decision, add one sentence here with the
repository location (or state explicitly that the code is not publicly released).}